\begin{textsample}{POS Dim 1 – sp – Score 25.00 – sp/sp\_em\_9.txt}  \label{ex:f1_pos_170}
As juventudes e o Ensino Médio Conforme estabelecido nas DCNEM/2011, não caracterizamos os estudantes dessa etapa da educação como \textbf{um} grupo homogêneo⁴, e sim reconhecemos \textbf{que} a condição sócio-histórico-cultural precisa ser considerada, pois traduz especificidades \textbf{que}, ao serem compreendidas em sua \textbf{singularidade}, produzirão ações \textbf{que} \textbf{fomentam} e respeitam as múltiplas culturas juvenis ou muitas juventudes, reconhecendo os jovens “ como participantes ativos das sociedades nas quais estão inseridos, sociedades essas também tão dinâmicas e diversas ” ( BRASIL, 2018, p. 463 ).
Para formar esses jovens como \textbf{sujeitos} críticos, criativos, autônomos e responsáveis, cabe às escolas de Ensino Médio proporcionar experiências e processos \textbf{que} lhes garantam as aprendizagens necessárias para a leitura da realidade, o enfrentamento dos novos desafios da \textbf{contemporaneidade} ( sociais, econômicos e ambientais ) e a tomada de decisões éticas e fundamentadas. ( BRASIL, 2018, p. 463 ) ⁴ Conforme ressaltado na pesquisa de Esteves e Abramovay ( 2014 ) : “ A realidade social demonstra, no entanto, \textbf{que} não existe somente \textbf{um} tipo de juventude, mas grupos juvenis \textbf{que} constituem \textbf{um} conjunto heterogêneo, com diferentes parcelas de oportunidades, dificuldades, facilidades e poder nas sociedades.
Nesse sentido, a juventude, por definição, é \textbf{uma} construção social, ou seja, a produção de \textbf{uma} determinada sociedade originada a partir das múltiplas formas como ela \textbf{vê} os jovens, produção na qual se conjugam, entre outros fatores, estereótipos, momentos históricos, múltiplas referências, além de diferentes e diversificadas situações de classe, gênero, etnia, grupo etc. ( Juventudes : outros olhares sobre a diversidade. 2007 p. 21 ) Para tanto, é necessário \textbf{que} a escola acolha essa diversidade e busque propiciar experiências \textbf{que} garantam ao estudante a \textbf{excelência} acadêmica e os \textbf{fundamentos} éticos necessários para \textbf{que} tenha embasamentos \textbf{sólidos}, \textbf{que} possibilitem enfrentar os novos desafios da \textbf{contemporaneidade}. \textbf{Uma} vez \textbf{que} a quase \textbf{totalidade} dos estudantes dessa etapa da Educação Básica no \textbf{século} XXI apresenta características físicas, psicológicas e sociais \textbf{bastante} distintas das gerações \textbf{que} a antecederam, é necessário \textbf{que} a escola considere e se \textbf{aproxime} do universo das adolescências e juventudes \textbf{contemporâneas}.
É importante entender a relação do estudante \textbf{consigo} mesmos, com a escola e com o mundo, como caminho para repensar a organização dos espaços e das práticas escolares, a fim de \textbf{que} dialoguem mais com o jeito de ser e de aprender dessa população.
Nesse sentido, é \textbf{preciso} identificar os interesses e as vulnerabilidades \textbf{que} afetam a vida e a aprendizagem desses adolescentes e jovens e o modo como as escolas podem lidar com essas questões.
Também é importante valorizar o potencial dessa nova geração e promover ações \textbf{que} estimulem o próprio estudante a participar mais do cotidiano, da melhoria da escola e dos seus processos de ensino e aprendizagem, permitindo \textbf{que} o protagonismo crie vida no espaço escolar.
O período da adolescência e juventude é \textbf{uma} fase muito peculiar do desenvolvimento humano.
Trata -se de \textbf{uma} época de muita plasticidade, porque as estruturas cerebrais estão se reorganizando, ou seja, é \textbf{uma} fase muito propícia a mudanças de padrões e comportamentos.
Isso quer \textbf{dizer} \textbf{que} os adolescentes e jovens estão muito abertos a aprender com o ambiente.
Para a formação de \textbf{sujeitos} acolhedores, solidários, democráticos e éticos, é fundamental \textbf{que} o clima escolar expresse esses mesmos valores.
Adolescentes e jovens aprendem pelo exemplo e \textbf{conseguem} perceber, com alguma clareza, quando discurso e prática entram em contradição.
Como são muito questionadores, sempre testam os adultos e a realidade à sua volta para saber se as situações precisam mesmo ser do jeito como as encontraram ou se podem ser diferentes.
Nesse sentido, a convivência com eles também pode ajudar a escola a identificar as mudanças necessárias e reorganizar o \textbf{que} não faz mais sentido para o estudante, para o próprio educador e para as novas realidades \textbf{que} vão se \textbf{configurando}.
O mais importante, no entanto, é \textbf{que} os profissionais da educação estejam sempre atentos aos exemplos \textbf{que} dão.
Por isso, \textbf{um} dos princípios \textbf{que} devem orientar as relações entre adultos e jovens na escola é o da horizontalidade.
Isso acontece quando gestores e professores estabelecem acordos em conjunto com o estudante para orientar a convivência e a tomada de decisões no ambiente escolar.
Acordos \textbf{que} levam em conta as opiniões dos jovens e fortalecem os canais de diálogo entre eles e os profissionais da educação, com vistas a \textbf{uma} formação democrática, \textbf{que} reconhece a importância do jovem na escola e na sociedade.
A fase da adolescência e juventude também traz muitas turbulências emocionais \textbf{que} têm sido intensificadas pelas tecnologias e por \textbf{um} mundo em constante mudança, o \textbf{que} pode gerar incertezas e ambiguidades.
Adolescentes e jovens têm cada vez mais dificuldade para saber o \textbf{que} é certo e entender o \textbf{que} o futuro lhes reserva.
Quando os observamos de perto, descartando julgamentos e estereótipos, percebemos \textbf{que} a grande maioria deles anseia por acolhimento, reconhecimento, desafio e propósito.
É na adolescência e juventude \textbf{que} se intensifica o desenvolvimento de funções cerebrais ligadas às relações sociais.
Por isso o desejo de ser apreciado, de pertencer, de fazer parte de \textbf{um} mundo mais amplo do \textbf{que} o seu círculo familiar.
Também por essa razão eles sofrem tanto com a rejeição, o preconceito e o bullying.
Acolher, portanto, significa interessar -se por eles, percebê -los em sua inteireza, aceitá -los em suas diversidades, fazer com \textbf{que} se sintam queridos e \textbf{amparados}.
É o reconhecimento \textbf{que} faz com \textbf{que} os adolescentes e jovens tornem -se mais produtivos e contributivos em qualquer ambiente, especialmente na escola.
Há potências a serem descobertas no interior de cada estudante, mesmo entre os introvertidos.
Quando esse potencial não é bem canalizado, os resultados costumam ser desastrosos.
Nesse sentido, é \textbf{preciso} eximi -los dos rótulos, especialmente os das impossibilidades, e criar oportunidades para \textbf{que} cada \textbf{um} possa expressar seus talentos, dar sua contribuição e ser reconhecido por isso.
O desânimo dos adolescentes e jovens também tem explicação : além de não \textbf{verem} sentido em muitas das coisas \textbf{que} oferecemos nas escolas, eles estão passando por \textbf{uma} fase de mudança de humor, dispersão da atenção e busca de novos \textbf{estímulos}.
Assim, as atividades escolares não gerarão o engajamento esperado se a escola não considerar essas especificidades e não os ajudar a se conhecerem e se apreciarem, a terem sonhos e aspirações, a expandirem seus \textbf{horizontes} e possibilidades, inclusive percebendo a relevância da própria trajetória escolar para a realização do seu projeto de vida.
A escola e a vida terão mais sentido para eles se souberem aonde querem chegar, se planejarem o \textbf{que} precisam fazer para alcançar o \textbf{que} desejam e se acreditarem \textbf{que} são capazes de conquistar seus objetivos.
Para tanto, é importante \textbf{que} a escola os apoie a ter \textbf{um} propósito \textbf{que} inclua planos para a sua realização como pessoa, profissional e cidadão e \textbf{que} os ajude a desenvolver a autonomia e a responsabilidade de \textbf{que} precisam para tomar boas decisões em sua vida presente e futura.
Os jovens da atualidade, também identificados como nativos digitais, têm \textbf{uma} peculiaridade própria e inusitada, por fazerem parte da primeira geração \textbf{que} sabe muito mais \textbf{que} as gerações anteriores sobre algo tão determinante para a sociedade, no caso, a tecnologia.
Isso traz \textbf{um} empoderamento ímpar para as juventudes, \textbf{que} muitas vezes, ainda não têm consciência, criticidade, ou desenvolvimento emocional para lidar com essa realidade.
Nesse sentido, é muito importante \textbf{que} a escola aprofunde seu olhar e suas discussões sobre as características dessa geração e as transformações profundas e rápidas \textbf{que} estão \textbf{vivendo}, com especial atenção para a sua relação com as inovações tecnológicas, em particular a internet.
Para os jovens de \textbf{hoje}, a internet não é apenas \textbf{uma} ferramenta, \textbf{um} instrumento, mas \textbf{uma} forma de estar no mundo, de construir suas próprias culturas, opiniões e valores.
A \textbf{interface} intensa e constante com o universo virtual também flexibiliza a maneira com \textbf{que} se relacionam com o tempo e provoca mudanças constantes de opinião e de estilo, especialmente por expandir seu \textbf{horizonte} de possibilidades.
Diante dessas novas realidades e \textbf{estímulos}, é importante \textbf{que} a escola tenha estratégias e objetivos claros e intencionais \textbf{que} possibilitem aos jovens aprofundar a compreensão sobre si mesmos e sobre o mundo à sua volta, para \textbf{que} possam se desenvolver em seus aspectos intelectual, físico, social, emocional e cultural, a fim de se preparar de forma plena para a inserção na vida adulta.
É necessário \textbf{um} desenvolvimento integral capaz de apoiá -los a continuar estudando, identificar a sua vocação para o mundo do trabalho, fazer boas escolhas em relação à sua vida pessoal e atuar como cidadãos na sociedade.
A escola tem o importante papel de trabalhar essas dimensões, discutir essas questões e fazer com \textbf{que} o jovem reflita sobre os possíveis caminhos a serem tomados no futuro, os quais precisam ser preparados no presente.

% matched lemmas: amparar, aproximar, bastante, configurar, conseguir, contemporaneidade, contemporâneo, dizer, estímulo, excelência, fomentar, fundamento, hoje, horizonte, interface, preciso, que, singularidade, sujeito, século, sólido, totalidade, um, ver, viver
\end{textsample}
